\chapter{Analysis and Discussion}
\label{ch:disc}

\section{Interpretation of the Results}
\label{sec:interpretation}

\TODO{This section interprets rather than restates Chapter~\ref{ch:impl}. Write
it after the numbers are in, structured as: what the key findings mean for each
hypothesis; which results support them and which do not; how the findings
compare with prior work; explanation of discrepancies; unexpected findings.}

The three hypotheses of Section~\ref{sec:rq} address separable mechanisms, and
the argument for each is independent of the others.

\paragraph{H1 --- formulation.} The claim is that posing retargeting as one
weighted functional dominates the sequential pipeline on both fidelity and
smoothness simultaneously. The distinction matters because the sequential
pipeline can be made arbitrarily smooth by increasing filter strength, at the
cost of fidelity; the interesting question is therefore not whether either
metric can be improved in isolation but whether the trade-off frontier itself
moves. \ph{[Insert result.]} If the proposed method improves one metric while
degrading the other, the correct conclusion is that the frontier has not moved
and the improvement is a re-parameterisation of the same trade-off.

\paragraph{H2 --- confidence weighting.} The mechanism proposed in
Section~\ref{sec:retarget} predicts a specific signature: the benefit of
weighting should concentrate on frames where perception is degraded, because
that is where uniform weighting drags the solution toward a bad observation.
A uniform improvement across occlusion strata would indicate that the weighting
is acting as a global regulariser rather than through the proposed mechanism.
\ph{[Insert result.]}

\paragraph{H3 --- replay screening.} The claim is about dependability rather
than accuracy: screening should raise the executable fraction of the exported
dataset. \ph{[Insert result.]} The stronger claim---that screening improves
downstream policy success---requires the paired training run noted in
Section~\ref{sec:res-policy} and should not be asserted on the executability
metric alone.

\section{Comparison with Prior Work}
\label{sec:comparison}

Direct numerical comparison with the systems reviewed in
Chapter~\ref{ch:lit} is not meaningful, since none is evaluated on the same
task suite, embodiment or sensing configuration. What can be compared is
qualitative behaviour and the class of guarantees provided.

Relative to OKAMI \cite{okami}, the present system discards the parametric
body-and-hand reconstruction and the two-stage warp-then-solve structure. The
expected consequence is a lower perception error floor---time-of-flight depth
in place of a mesh fit---at the cost of being restricted to parallel-jaw
targets. Relative to ORION \cite{orion}, the present system retains the dense
wrist trajectory that ORION discards, which should preserve velocity-profile
information at the cost of a dependency on reliable hand observation. Relative
to the online teleoperation systems \cite{anyteleop, opentv, ace}, the
formulation is the same but the operating regime is not: the offline setting
permits whole-trajectory optimisation with a well-defined acceleration term at
every interior timestep and admits a physical validation stage, neither of
which an interactive-rate solver can perform.

The measured cross-camera registration error
(Table~\ref{tab:calib}) should be compared against the $7.4$--$21.4$~mm range
reported by Romeo et al.\ \cite{romeo}. \TODO{State whether the measured value
falls in that range and, if not, diagnose the cause rather than reporting the
discrepancy without explanation.}

\section{Error Budget}
\label{sec:errbudget}

A central methodological commitment of this work is that residual error is
decomposed rather than reported as a single aggregate figure. Three
contributions are independent in origin and demand different remedies;
conflating them would make the aggregate uninformative about what to improve.

\begin{enumerate}[leftmargin=2em]
  \item \emph{Landmark detection noise.} Error in the image-plane hand
        landmarks, propagated through depth lifting. This contribution is
        reduced by confidence weighting \eqref{eq:weights} and by multi-view
        fusion, and it is characterised by the wrist-position repeatability
        measurement of Table~\ref{tab:calib}. It is not reduced by better
        calibration.
  \item \emph{Calibration error.} Error in the camera-to-workspace and
        robot-to-workspace transforms \eqref{eq:extrinsic}. This contribution
        is systematic within a session, appears as a constant offset rather
        than as noise, and is characterised by the reprojection and
        cross-camera registration measurements. It is not reduced by
        multi-view fusion, since all views inherit it.
  \item \emph{Representation dependency.} Error introduced by the choice of
        reference frame, specifically the sensitivity of world-frame
        trajectories to object displacement. This contribution is zero when
        the object is where it was during demonstration and grows with
        displacement; it is characterised by the comparison of
        Table~\ref{tab:repr} rather than by any calibration measurement.
\end{enumerate}

\begin{table}[htbp]
\caption{Decomposition of Residual Wrist-Pose Error by Source}
\label{tab:errbudget}
\centering
\small
\begin{tabular}{llcc}
\toprule
Source & Character & Magnitude (mm) & Reduced by \\
\midrule
Landmark detection & random  & \ph{X.X} & fusion, confidence weighting \\
Calibration        & systematic & \ph{X.X} & board size, view count \\
Representation     & conditional & \ph{X.X} & object-relative frame \\
\midrule
Combined           & --- & \ph{XX.X} & --- \\
\bottomrule
\end{tabular}
\end{table}

\TODO{Populate and discuss. If the combined figure differs materially from the
quadrature sum of the independent contributions, the sources are not
independent and the decomposition must be revised rather than presented as if
it held.}

\section{Limitations}
\label{sec:limitations}

The assumptions of Section~\ref{sec:assumptions} translate into the following
limitations, which are stated plainly rather than deferred to future work.

\begin{enumerate}[leftmargin=2em]
  \item \emph{Gripper abstraction.} A binary gripper state cannot represent
        in-hand manipulation, regrasping or force-modulated grasps. Tasks
        requiring these are outside the demonstrated capability, and extending
        to dexterous hands would require reinstating the finger-level
        representation that Section~\ref{sec:lit-hand} rejects.
  \item \emph{Rigid-object assumption.} The object-relative transform
        \eqref{eq:objrel} is undefined for deformable objects. Cloth, cable and
        granular manipulation are excluded.
  \item \emph{Dependence on object pose estimation.} The object-relative
        representation inherits the error of the object pose estimator. Where
        that estimator fails---textureless, symmetric or occluded objects---the
        representation degrades to worse than the world-frame alternative it is
        meant to improve upon.
  \item \emph{Kinematic transfer only.} No force or contact model is estimated.
        Success on tasks with a force-regulation component is not claimed.
  \item \emph{Single demonstrator, constrained task class.} The evaluation uses
        \ph{one} demonstrator and \ph{three} task types. Generalisation across
        demonstrators, hand morphologies and task families is not established.
  \item \emph{Sensing floor.} Cross-camera agreement at the centimetre level
        \cite{romeo} bounds achievable accuracy. Tasks with tolerances below
        that bound cannot be addressed by this sensing configuration
        regardless of the retargeting formulation.
  \item \emph{Cross-family synchronisation.} Alignment between the Kinect pair
        and the RealSense is software-based and therefore millisecond-scale;
        fast motion may exhibit residual misalignment between those streams.
\end{enumerate}

\section{Threats to Validity}
\label{sec:validity}

\paragraph{Internal validity.} Metrics were fixed before the final runs
(Section~\ref{sec:metrics}) to prevent selection of criteria after observing
outcomes. Baselines B1--B5 differ from the proposed method in exactly one
component each, so attribution of any observed difference is unambiguous. The
principal residual threat is that the sequential baseline's filter parameters
are themselves tunable; an unfavourable choice would make the comparison
misleading. \TODO{State how B1 filter parameters were selected---ideally by
tuning them to optimise the baseline's own fidelity, not by adopting a default.}

\paragraph{External validity.} The constrained task class was chosen for
measurability, which necessarily limits generalisation. Results should not be
read as claims about manipulation in general.

\paragraph{Construct validity.} Spectral arc length measures smoothness of the
velocity profile, which is a proxy for the property actually of interest---that
a policy can imitate the trajectory. The proxy is standard but imperfect, and
the downstream policy results of Section~\ref{sec:res-policy} are the more
direct evidence.

\section{Implications}
\label{sec:implications}

Three implications extend beyond the specific system.

First, the transferability of the online retargeting formulation to the offline
regime suggests that the separation between the teleoperation and
video-imitation literatures is a historical accident rather than a technical
necessity. The offline setting is strictly less constrained---no causality
requirement, no latency budget---and therefore able to use the same formulation
more thoroughly than the setting for which it was designed.

Second, moving curation upstream of export, so that infeasible trajectories are
removed before they enter the dataset rather than filtered after they have
degraded a policy, is applicable to any pipeline that synthesises rather than
records actions. The cost is one physical execution per episode; the benefit is
that the exported dataset contains proprioception the robot actually attained.

Third, the error decomposition of Section~\ref{sec:errbudget} is a
transferable methodological point. Reporting a single aggregate accuracy figure
for a multi-stage perception pipeline obscures which stage limits performance,
and consequently which engineering effort would pay.
